\documentclass[10pt,a4paper]{amsart}
\usepackage[margin=19mm]{geometry}
\usepackage{amsmath,amssymb,mathtools}
\usepackage[scr=rsfs,cal=euler]{mathalfa}
\usepackage{xfrac}
\usepackage[unicode,hidelinks,bookmarks=false]{hyperref}
\newcommand{\ie}{i.e.\ }
\newcommand{\cf}{cf.\ }
\newcommand{\resp}{respectively\ }
\newcommand{\Subsection}{\subsection}
\providecommand{\nobreakdash}{\nobreak\hbox{-}}
\setlength{\emergencystretch}{2em}
\multlinegap0pt
\usepackage{bm}
\usepackage{bbm}
\usepackage{dsfont}
\usepackage{xspace}
\usepackage{cancel}
\usepackage{mathtools}
\usepackage{amsthm}
\makeatletter
\@ifundefined{theorem}{\newtheorem{theorem}{Theorem}}{}
\makeatother
\newcommand{\LPF}{\mathsf{LuckyPF}}  
\newcommand{\bsy}{\boldsymbol}
\newcommand{\out}{\mathcal{O}}
\title[Enumerating Vector Parking Functions and their Outcomes Base]{Enumerating Vector Parking Functions and their Outcomes Based on Specified Lucky Cars}
\author{Melanie Ferreri, Pamela E. Harris, Lucy Martinez, Eric Swartz}
\date{Source: arXiv:2508.13917v2 (CC BY 4.0)}
\begin{document}
\maketitle

\noindent\emph{Source and licence.} Enumerating Vector Parking Functions and their Outcomes Based on Specified Lucky Cars. Version arXiv:2508.13917v2 (CC BY 4.0), distributed under the Creative Commons Attribution 4.0 International licence, as recorded in the arXiv licence metadata for this exact version and shown on the abstract page https://arxiv.org/abs/2508.13917v2. Full rights notice: licenses/arxiv-2508.13917-CC-BY-4.0.txt.\par\smallskip

\noindent\textit{Editorial scope.} Counted unit: the Theorem whose source label is \texttt{thm:countforupfs} in the article (source0.tex lines 1154--1166, source label \texttt{thm:countforupfs}), delivered together with the article's own complete original proof of it (source0.tex lines 1168--1221, one \texttt{proof} environment of 54 lines). No mathematics is written, repaired or re-derived: statement and proof are the article's own text line for line. Required context retained uncounted: none. Every definition, display and label of those lines is retained unchanged. Omitted from this package and not redistributed: 1--1153, 1167--1167, 1222--1338. Nothing inside a retained excerpt is omitted or altered: every retained body is delivered exactly as the article's lines are written, so no omission receipt of any kind accompanies this package. Citations: the retained excerpts carry no citation command at all, so no bibliography entry of the article is transcribed and no reference is redistributed. Reference adaptation: none was needed and none was made; every reference inside the delivered unit resolves inside it, so no unresolved reference or citation is produced. Printed slips kept literal: no printed word, symbol, hypothesis or number of the counted statement or of its proof was altered. Layout and class: the article's own class is \texttt{amsart} with packages including inputenc, xcolor, fullpage, hyperref, cleveref, enumitem, tikz, amsmath, amssymb, color, mathtools, graphicx, geometry, amsthm; the wrapper is the lane's pinned \texttt{amsart} editorial wrapper at 10pt on A4, which restates the theorem-like environments the retained text needs and adds the article's own macro definitions that the retained text needs (\texttt{\textbackslash LPF}, \texttt{\textbackslash bsy}, \texttt{\textbackslash out}), with extra wrapper packages (bm, bbm, dsfont, xspace, cancel, mathtools). The delivered numbering is the unit's own. Assets: no retained span contains a figure environment, a graphics-inclusion command, a PDF-page inclusion, a table or a TikZ picture; no image file is copied into this package, the package declares no asset dependency and no image file is redistributed with it. Licence: version arXiv:2508.13917v2 of this article is distributed under the Creative Commons Attribution 4.0 International licence, as recorded in the arXiv licence metadata for this exact version and shown on the abstract page https://arxiv.org/abs/2508.13917v2. A full rights notice is delivered at licenses/arxiv-2508.13917-CC-BY-4.0.txt. Characters: every retained excerpt is pure ASCII, so the delivered engine, XeLaTeX with its default Latin Modern text font, reports no missing character.
\begin{theorem}\label{thm:countforupfs}
Let $\bsy{u}=(\underbrace{u_1,u_1,\ldots, u_1}_{m_1}, \underbrace{u_2,u_2,\ldots, u_2}_{m_2}, \ldots, \underbrace{u_r,u_r,\ldots ,u_r}_{m_r})$ where for each $u_j\in \bsy{u}$, $u_j$ has multiplicity $m_j$ for all $j\in [r]$. Then the number of $\bsy{u}$-parking functions with $k$ lucky cars is
\[ 
\left|\bigcup_{|I|=k} \LPF_{\bsy{u}}(I)\right| = 
\sum_{{\substack{\kappa = (k_1,k_2,\ldots, k_r), \\
\sum_{i=1}^{r}k_i=k,\\
0\leq k_i\leq m_i}}}
\,\sum_{{B = (B_1,B_2,\ldots, B_r) \in \out_{\bsy u}}}\binom{m_1}{k_1}(u_1-1)^{m_1-k_1} S(B ; \kappa),
\]
where
\[S(B;\kappa) \coloneq \sum_{{\substack{(L_2,L_3,\ldots,L_r), \\ L_d\subseteq B_d, |L_d|=k_d}}}\prod_{j=2}^r \prod_{t=1}^{j} (u_{j}-u_{t-1}-1)^{\ell_{j,t}-\ell_{j,t-1}}.\]

\end{theorem}

\begin{proof}
Recall that in this subsection, $B_i$ is no longer the set of cars that parked in spot $i$, but given the reindexing it is now is the $i$th nonempty set that appears in the outcome. 
By assumption, we have $r$ spots with capacities given by $m_1,\ldots, m_r$. Fix an outcome $B =(B_1,B_2,\ldots, B_r)\in\out_{\bsy{u}}$, and suppose that there are $k$ lucky cars. Let $\kappa = (k_1,k_2,\ldots, k_r)$ be a weak composition of $k$, such that the $i$th spot contains $k_i$ lucky cars. Then we have that $0 \leq k_i \leq m_i$ for all $i$.
We count the number of parking functions corresponding to this outcome.

The first available spot is $u_1$.
We have that $B_1$ is the set of cars which parked in spot $u_1$, and $k_1$ of them are lucky. 
First we choose which cars are lucky; there are $ \binom{m_1}{k_1}$ options. For these cars, their preference is fixed to be $u_1$. 

Then there are $m_1-k_1$ unlucky cars remaining in $B_1$. We count the possible preferences for the unlucky cars in this set. 
For each such car, their preference could have been any spot earlier than $u_1$, so they each have $(u_1-1)$ options for their preference. Hence, the number of ways to count the preferences of cars in the set $B_1$ is given by 
$\binom{m_1}{k_1}(u_1-1)^{m_1-k_1}$.

We now consider the possible preferences for the cars in spots $B_2,\ldots, B_r$. Let $S(B,\kappa)$ denote the number of possible preferences among the cars parking in the remaining $r-1$ available spots.

Fix $u_j$ to be one of the spots after $u_1$, and let $L_j$ be the set of lucky cars in spot $u_j$ (so $L_j$ is a subset of $B_j$ consisting of $k_j$ cars). 
For the cars in $L_j$, their preference is fixed to be $u_j$ (the spot where they parked). 
We consider sets of  possible preferences for the remaining $m_j-k_j$ unlucky cars.



Fix a spot $u_t\leq u_{j-1}$ and define $u_0 =0$. 
We count the number of unlucky cars $c$ in $B_j$ whose set of possible preferences contains any spot between $u_{t-1}$ and $u_j$, not including $u_{t-1}$ (or $u_j$, since $c$ is an unlucky car). When such a car $c$ went to park, its preference was after spot $u_{t-1}$, so it did not check this spot, and all spots from $u_t$ to $u_{j-1}$ were full. However, spot $u_{t-1}$ would not have been full since if it was, then car $c$ could have preferred spot $u_{t-1}$ as well. 

The number of such cars is given by $\ell_{j,t}-\ell_{j,t-1}$. These cars could have preferred any spot in the set $\{u_{t-1}+1, ,u_{t-1}+2,\ldots, u_{j}-1\}$, and in particular not any spot $u_{t-1}$ or earlier, to end up parking in spot $u_j$.

For each such car, there are {$u_{j}-1 - (u_{t-1} +1 ) +1 = u_{j} - u_{t-1} -1 $} options for its preference. 
Then the number of possible preferences among all cars which could have preferred any spot $u_t \leq u_{j-1}$ is given by $\prod_{t=1}^{j-1} (u_{j}-u_{t-1}-1)^{\ell_{j,t}-\ell_{j,{t-1}}}$.

The unlucky cars that remain to be accounted for are those that preferred a spot between $u_{j-1}$ and $u_j$, but were unable to park there since it was unavailable.
There are {$u_j-u_{j-1}-1$} unavailable spots between $u_{j-1}$ and $u_j$. To avoid double-counting cars that were accounted for in the previous product, we only consider cars such that when they parked, spot $u_{j-1}$ was available; in other words spot $u_{j-1}$ was not full yet. This guarantees that their preference could not have been $u_{j-1}$ or earlier (or else they would have parked there instead).


The number of unlucky cars such that spot {$u_{j-1}$} was {available} when they parked is given by $m_j-k_j-\ell_{j,j-1}$. 
It follows that $(u_j-u_{j-1}-1)^{m_j-k_j-\ell_{j,j-1}}$ is the number of possible preferences among unlucky cars preferring a spot after $u_{j-1}$ and parking in spot $u_j$.
In total we have, for a fixed spot $u_j$, $$ \left(\prod_{t=1}^{j-1} (u_{j}-u_{t-1}-1)^{\ell_{j,t}-\ell_{j,t-1}}\right)\left(u_j-u_{j-1}-1\right)^{m_j-k_j-\ell_{j,j-1}}$$ possible preferences for the cars in the set $B_j$.
With the convention that $\ell_{j,j} = m_j - k_j$, this product simplifies to
\[ \prod_{t=1}^{j} (u_{j}-u_{t-1}-1)^{\ell_{j,t}-\ell_{j,t-1}}.\] 
Taking the product over all $j=2,\ldots, r$, we have for
a given outcome $B = (B_1,B_2,\ldots, B_r)$,  with lucky cars in each of the spots given by $(L_2,\ldots, L_r)$, a total of
\[ \prod_{j=2}^r \prod_{t=1}^{j} (u_{j}-u_{t-1}-1)^{\ell_{j,t}-\ell_{j,t-1}}\] possible preference lists for the cars in spots $2$ through $r$.
Summing over all possible distributions of lucky cars $(L_2,\ldots, L_r)$, we have 
\[S(B;\kappa) = \sum_{\substack{(L_2,L_3,\ldots,L_r), \\L_d\subseteq B_d, |L_d|=k_d}}\prod_{j=2}^r \prod_{t=1}^{j} (u_{j}-u_{t-1}-1)^{\ell_{j,t}-\ell_{j,t-1}}\]
total options for the preferences of cars parking in the last $r-1$ spots. Then for a fixed outcome $B = (B_1,B_2,\ldots, B_r) \in \out(\bsy u)$, the total number of corresponding parking functions is
\[
\binom{m_1}{k_1}(u_1-1)^{m_1-k_1} S(B ; \kappa).
\]
Since each parking function corresponds to one outcome, we may enumerate all parking functions with $k$ lucky cars by summing over all possible outcomes with $k$ lucky cars. Then the number of parking functions of the desired form is 

\[
\sum_{\substack{\kappa = (k_1,\ldots, k_r),\\0\leq k_i\leq m_i,\\\sum_{i=1}^{r}k_i=k}}
\,\sum_{{B = (B_1,B_2,\ldots, B_r) \in \out(\bsy u)}}\binom{m_1}{k_1}(u_1-1)^{m_1-k_1} S(B ; \kappa). \qedhere
\]
\end{proof}

\end{document}
